% Chapter Template

% Main chapter title
%\chapter[toc version]{doc version}
\chapter{Conclusion}

% Short version of the title for the header
%\chaptermark{version for header}

% Chapter Label
% For referencing this chapter elsewhere, use \ref{ChapterTemplate}
\label{Conclusion}

% Write text in here
% Use \subsection and \subsubsection to organize text


\quad This thesis set out to design, implement, and evaluate a secure livestreaming framework capable of sustaining reliable delivery in dense wireless environments while reducing infrastructure load through fog offloading. The system was motivated by the growing demand for live video in high-density scenarios, where traditional multicast over Wi-Fi suffers from poor reliability and conventional solutions demand expensive infrastructure upgrades, the proposed architecture combined protocol-level redundancy, adaptive retransmission scheduling, peer-assisted recovery, and a security model tailored for untrusted fog nodes and opportunistic device-to-device links. 

\quad The implementation, described in \textbf{\Cref{Design and Implementation}}, integrated these mechanisms into a modular pipeline built in safe Rust, deployed through containers for reproducibility, and extended with an Android frontend leveraging Wi-Fi Aware mesh connectivity. Evaluation in \textbf{\Cref{Evaluation}} demonstrated that static FEC, although simple, consumed large amounts of bandwidth without proportional reliability gains, while targeted retransmissions halved missed deadlines at the same overhead cost, reinforcement learning showed potential for adaptive parity selection but in practice generated excessive redundancy under the chosen reward structure. The most decisive result came from real-world tests: fog offloading preserved baseline overhead yet reduced deadline misses by almost 80\%, validating the design goal of exploiting local peer recovery to improve packet delivery works without burdening the access point.

\quad Together, these results confirm that the framework secures the stream against adversarial peers, maintains efficient use of infrastructure bandwidth, and leverages fog resources to enhance playback continuity. By combining retransmission scheduling with mesh-based offloading, the system provides a practical path for scalable deployment in environments where reliability and efficiency are critical, such as stadiums, concerts, and large public events. This thesis contributes a complete architecture and empirical validation for secure livestreaming with fog offloading and establishes a foundation for future research on scalable, secure, and efficient video delivery at the network edge.


\section{Future Work}
\label{sec:future work}

\quad Future work should investigate extending the Android implementation with advanced peer-to-peer frameworks to broaden recovery strategies and increase system resilience. One promising direction is the integration of the \texttt{libp2p framework} ~\cite{libp2p}, which was partially explored in a separate development branch. The approach enabled compatibility by creating a virtual TAN interface on Android, routing all Wi-Fi Aware traffic through it, and broadcasting to peers so that each device obtained an IPv6 address linkable to a \texttt{libp2p multiaddress}. This integration allowed the system to experiment with multipath recovery across heterogeneous interfaces, such as combining Wi-Fi Aware with mobile data, and opened the possibility of richer peer-to-peer topologies with routing and other advanced features.  

\quad However, further development is required to make this approach practical, as Android’s TAN interface does not support multicast and is therefore incompatible with the mDNS peer discovery mechanism used in \texttt{libp2p}. A possible workaround is to introduce a bootstrap node to provide discovery services, but this would shift the topology from decentralized to centralized, conflicting with the requirements defined in \textbf{\Cref{sec:requirements}}. Future work should therefore focus on addressing this limitation, either by extending \texttt{libp2p} to support IPv6 link-local addressing used in Wi-Fi Aware or by designing an alternative decentralized discovery mechanism compatible with TAN-based routing. At the same time, reinforcement learning can be refined with latency-aware reward functions to prioritize deadline-sensitive recovery, multicast coding can be optimized to increase coding gain and reduce reliance on unicast retransmissions, and fog offloading must be validated on heterogeneous mobile devices and larger peer groups to confirm scalability and robustness. Within this broader scope, adopting \texttt{libp2p} offers a path toward multipath transport and richer peer-to-peer topologies while reducing the complexity and limitations of the custom mesh implementation described in \textbf{\Cref{sec:mesh networking}}. Advancing in these directions will enable the design of resilient peer-to-peer systems capable of sustaining reliable delivery in dense wireless environments without sacrificing decentralization.